% ============================================================
% KẾT LUẬN (không đánh số chương, ngang hàng với Tài liệu tham khảo)
% ============================================================
\chapter*{Kết luận}
\addcontentsline{toc}{chapter}{Kết luận}
\label{chap:ket-luan}

\noindent\textbf{Tóm tắt}
\medskip

Tiểu luận đã trình bày cơ sở lý thuyết của Thuật toán Di truyền
và mở rộng của nó là Lập trình Di truyền, sau đó minh họa cả hai
phương pháp qua bốn thực nghiệm cụ thể. Chương~\ref{chap:co-so-ly-thuyet}
xây dựng khung lý thuyết theo trình tự: tổng quan bài toán tối ưu và
họ thuật toán tiến hóa; cơ sở khoa học của GA với ba cách mã hóa cá
thể (nhị phân, số thực, hoán vị); các thành phần cốt lõi của một vòng
lặp tiến hóa (hàm thích nghi, chọn lọc, lai ghép, đột biến, elitism);
quy trình thuật toán và điều kiện dừng; và cuối cùng là GP với mã hóa
dạng cây biểu thức.

Trên nền tảng đó, Chương~\ref{chap:ung-dung} áp dụng GA cho ba nhóm
bài toán tối ưu số — không ràng buộc trên năm hàm benchmark kinh điển,
có ràng buộc trên năm bài toán chuẩn CEC 2006, và bài toán người du
lịch kết hợp 2-opt local search theo mô hình memetic algorithm — đối
chiếu trực tiếp với các mốc nghiệm tối ưu đã công bố. Ở cả ba nhóm,
GA (khi cần thiết được hỗ trợ thêm bởi tìm kiếm cục bộ) đều đạt hoặc
tiệm cận nghiệm tối ưu đã biết. Riêng ứng dụng thứ tư dùng GP cho bài
toán hồi quy symbolic, dự đoán nhiệt độ điểm sương từ nhiệt độ không
khí và độ ẩm tương đối trên dữ liệu thời tiết thực tế, đạt
$R^2 = 0.9967$ trên tập kiểm tra và tìm ra một biểu thức tường minh
phản ánh đúng cấu trúc phi tuyến (logarit, phép chia lồng nhau) của
quan hệ vật lý thật (công thức Magnus), thay vì chỉ đưa ra một mô hình
dự đoán dạng hộp đen.

\bigskip
\noindent\textbf{Hạn chế}
\medskip

Bên cạnh những kết quả đạt được, tiểu luận vẫn còn một số hạn chế cần
được xem xét. Thứ nhất, tính linh hoạt của GA/GP đi kèm với việc không
khai thác trước các đặc trưng cấu trúc của bài toán, do đó hiệu quả tìm
kiếm có thể chưa cao trong một số trường hợp. Đối với các bài toán tối
ưu có ràng buộc, GA thuần túy có thể gặp khó khăn trong việc tìm kiếm
các nghiệm nằm sát biên của miền khả thi. Kết quả thực nghiệm trên các
bài toán CEC 2006 cho thấy việc bổ sung cơ chế tìm kiếm cục bộ có thể
giúp thuật toán tiếp cận nghiệm tốt hơn. Tương tự, đối với bài toán
người du lịch (TSP) có quy mô lớn, việc kết hợp GA với các phép cải
thiện nghiệm như 2-opt giúp quá trình hội tụ nhanh hơn và đưa nghiệm
thu được đến gần nghiệm tối ưu.

Thứ hai, trong bài toán hồi quy biểu thức bằng GP, không gian biểu diễn
hằng số của các nút lá còn bị giới hạn. Cụ thể, engine GP được xây dựng
trong tiểu luận chỉ cho phép các hằng số nằm trong khoảng $[-2,2]$.
Điều này gây khó khăn khi biểu diễn chính xác các hằng số có độ lớn
lớn hơn, chẳng hạn các hằng số $a, b$ xuất hiện trong công thức Magnus.
Để biểu diễn những giá trị này, GP phải kết hợp nhiều nút và phép toán
trung gian, từ đó làm tăng kích thước cây và tiêu tốn ngân sách độ sâu
vốn cần thiết cho việc xây dựng cấu trúc chính của biểu thức. Đây cũng
là một trong những nguyên nhân khiến GP chưa thể khôi phục chính xác
tuyệt đối công thức Magnus.

Thứ ba, phạm vi thực nghiệm của tiểu luận còn tương đối hạn chế. Các
thí nghiệm hiện mới được thực hiện với một thiết lập ngẫu nhiên cố
định cho mỗi cấu hình tham số, do đó chưa đánh giá đầy đủ mức độ ổn
định của thuật toán trước sự thay đổi của quá trình khởi tạo. Bên cạnh
đó, quy mô của các bài toán được khảo sát, chẳng hạn số chiều trong các
bài toán tối ưu có ràng buộc, số địa điểm trong TSP hay số lượng và
đặc trưng của biến đầu vào trong hồi quy biểu thức, vẫn còn tương đối
nhỏ so với nhiều bài toán thực tế phức tạp. Vì vậy, khả năng mở rộng và
tính ổn định của GA/GP trên các bài toán có quy mô lớn hơn vẫn cần
được đánh giá thêm.

\bigskip
\noindent\textbf{Hướng phát triển}
\medskip

Từ những hạn chế đã nêu, tiểu luận đề xuất một số hướng phát triển cho
các nghiên cứu tiếp theo:

\begin{itemize}
    \item Đối với hạn chế về hiệu quả tìm kiếm, có thể nghiên cứu việc
    kết hợp GA với các phương pháp tìm kiếm cục bộ phù hợp với từng
    bài toán. Bên cạnh 2-opt cho TSP, có thể áp dụng các phương pháp
    tối ưu dựa trên gradient đối với những bài toán có hàm mục tiêu và
    ràng buộc khả vi. Đồng thời, có thể sử dụng cơ chế tham số thích
    ứng để tự điều chỉnh tỷ lệ lai ghép, tỷ lệ đột biến và kích thước
    quần thể theo tiến trình hội tụ.

    \item Đối với hạn chế về khả năng biểu diễn hằng số trong GP, có
    thể bổ sung bước tối ưu riêng cho các hằng số tại nút lá sau khi
    cấu trúc cây được xác định. Các phương pháp như least-squares hoặc
    gradient descent có thể được sử dụng để tinh chỉnh các hằng số,
    thay vì chỉ dựa vào đột biến ngẫu nhiên trong một khoảng giá trị
    cố định.

    \item Đối với hạn chế về phạm vi thực nghiệm, nên thực hiện mỗi
    cấu hình nhiều lần với các quá trình khởi tạo ngẫu nhiên khác
    nhau và báo cáo giá trị trung bình cùng độ lệch chuẩn. Đồng thời,
    có thể tăng quy mô bài toán, chẳng hạn số chiều, số địa điểm trong
    TSP và số lượng đặc trưng đầu vào, nhằm đánh giá tốt hơn tính ổn
    định và khả năng mở rộng của GA/GP.
\end{itemize}
